% The derivation that exists. Sign derived from Einstein+Λ on SdS.
% Local Clausius ⇒ Einstein+Λ (Jacobson 1995). Vacuum is dS.
% Not FGHMV in cosmology. Not a prize theorem.
\documentclass[11pt]{article}
\usepackage[margin=1.05in]{geometry}
\usepackage{amsmath,amssymb,amsthm}
\usepackage{enumitem,booktabs}
\usepackage[colorlinks=true,linkcolor=blue,citecolor=blue,urlcolor=blue]{hyperref}

\newtheorem{theorem}{Theorem}[section]
\newtheorem{lemma}[theorem]{Lemma}
\newtheorem*{admission}{Admission}
\newcommand{\pP}{\textbf{[P]}}

\title{\textbf{The Derived Sign:\\
Local Clausius Gives Einstein$+\Lambda$,\\
and Einstein$+\Lambda$ Gives the Cosmological Minus Sign}}
\author{Robert W.\ McGwier\\
\small CTO, Cohere Technology Group\\
\small GitHub: \texttt{n4hy}}
\date{15 August 2026 --- Version 1}

\begin{document}
\maketitle

\begin{abstract}
\noindent
This is the derivation that exists. It is not Faulkner--Guica--Hartman--
Myers--Van Raamsdonk on a net of de~Sitter balls, and it is not a
Nobel theorem.

Locally: the Unruh temperature is positive, $\delta Q=T\,dS$ on a
local Rindler horizon with $dS=\eta\,\delta A$ implies Einstein's
equation with an undetermined $\Lambda$ \cite{Jacobson1995}. That
algebra is \pP{} (Paper~18). The spinless vacuum with $\Lambda>0$
is de~Sitter (Paper~19). Torsion is not in the problem.

Globally, on the cosmological horizon of Schwarzschild--de~Sitter,
the same field equation implies
\begin{equation*}
T\,\mathrm{d}S+\mathrm{d}M=0.
\end{equation*}
At $M=0$, $\mathrm{d}r_c/\mathrm{d}M=-1$ and
$T\,\mathrm{d}S=-\mathrm{d}M$. The minus sign is \emph{derived}
from Einstein$+\Lambda$, not copied from AdS. \pP{} (implicit
differentiation, audited; finite-$M$ residual $<10^{-6}$).

Copying the Casini--Huerta--Myers plus sign onto the Gibbons--Hawking
horizon remains false (Paper~17). Using the minus sign on that
horizon is using Einstein$+\Lambda$. The two signs live on two
surfaces and are consistent.

What is not \pP{}: a holographic pair, a value of $\Lambda$,
FGHMV equivalence for arbitrary balls in de~Sitter.
\end{abstract}

\section{Local implication (cited, already isolated)}

Unruh: $T=\hbar\kappa/2\pi>0$ for a boost. Clausius on a local
Rindler horizon with $\theta=\sigma=0$ at the bifurcation surface,
together with $dS=\eta\,\delta A$ and $\nabla^a T_{ab}=0$, implies
\cite{Jacobson1995,McGwierXVIII}
\begin{equation}
G_{\mu\nu}+\Lambda g_{\mu\nu}
=
\frac{2\pi}{\hbar\eta}\,T_{\mu\nu},
\qquad
\Lambda\text{ undetermined}.
\label{eq:loc}
\end{equation}
The Unruh sign is plus. It is a theorem of QFT, not an AdS import.
The area law is an assumption. \pP{} as an implication under those
hypotheses.

\section{Vacuum}

If $T_{\mu\nu}=0$ and $\Lambda>0$, \eqref{eq:loc} is de~Sitter of
radius $\ell=\sqrt{3/\Lambda}$ (Paper~19). In the flat FLRW chart,
$a=e^{Ht}$ with $H^2=\Lambda/3$ saturates both Friedmann equations.
Torsion vanishes with the spin. \pP{}

\section{The cosmological sign, derived}

Schwarzschild--de~Sitter ($G=c=\hbar=1$)
\begin{equation}
f(r)=1-\frac{2M}{r}-\frac{r^2}{\ell^2},
\qquad
\Lambda=\frac{3}{\ell^2}.
\label{eq:sds}
\end{equation}
The cosmological root $r_c(M)$ of $f(r)=0$ satisfies $r_c(0)=\ell$.
Implicit differentiation,
\begin{equation}
\frac{\mathrm{d}r_c}{\mathrm{d}M}
=
-\frac{\partial_M f}{\partial_r f}
=
\frac{\ell^2 r}{M\ell^2-r^3},
\qquad
\left.\frac{\mathrm{d}r_c}{\mathrm{d}M}\right|_{M=0}=-1.
\label{eq:imp}
\end{equation}
Surface gravity $\kappa=\lvert f'(r_c)\rvert/2$ equals $1/\ell$ at
$M=0$, so $T=\kappa/2\pi=1/(2\pi\ell)$. Entropy $S=\pi r_c^2$ then
gives
\begin{equation}
T\,\mathrm{d}S
=
\frac{1}{2\pi\ell}\cdot 2\pi\ell\cdot\mathrm{d}r_c
=
-\mathrm{d}M
\qquad(M=0).
\label{eq:minus}
\end{equation}
The first law of the cosmological horizon is
$T\,\mathrm{d}S+\mathrm{d}M=0$. The auditor recovers
\eqref{eq:imp}--\eqref{eq:minus} identically in sympy. At finite
small $M$ the same combination is zero to $10^{-6}$. Without the
$\Lambda$ term there is no cosmological root. \pP{}

This is the Gibbons--Hawking minus sign
\cite{GibbonsHawking1977,Banihashemi2023}, now as an identity of
\eqref{eq:sds}, not a citation standing alone.

The physical-process reading is the same identity: positive Killing
energy crossing the past horizon focuses the generators
(Raychaudhuri plus Einstein), so the area at the bifurcation surface
falls \cite{Banihashemi2023}.

\section{Two surfaces, two signs, no contradiction}

\begin{center}
\begin{tabular}{@{}lll@{}}
\toprule
Surface & Sign & Source \\
\midrule
Local Rindler / CHM ball & $+\delta\langle T_{00}\rangle\Rightarrow+\delta S$
  & Unruh + Clausius \\
Cosmological horizon & $+\mathrm{d}M\Rightarrow-\mathrm{d}S$
  & Einstein$+\Lambda$ on SdS, \eqref{eq:minus} \\
\bottomrule
\end{tabular}
\end{center}

Putting the first sign on the second surface is false (Paper~17).
Putting the second sign on the second surface is Einstein$+\Lambda$.
The derivation does not copy AdS onto de~Sitter. It uses Unruh
locally and \eqref{eq:sds} globally.

\section{What this derivation is}

\begin{enumerate}[leftmargin=1.6em]
\item Local entanglement/Clausius with the Unruh plus sign
  $\Rightarrow$ Einstein$+\Lambda$ (hypotheses of \cite{Jacobson1995}).
\item Vacuum of that equation with $\Lambda>0$ $\Rightarrow$ de~Sitter.
\item The same equation on SdS $\Rightarrow$ cosmological
  $T\,\mathrm{d}S+\mathrm{d}M=0$.
\end{enumerate}
That is emergent bookkeeping yielding de~Sitter \emph{as the vacuum
of the metric sector it produces}, with the global sign derived.
It is not an FGHMV theorem about a cosmological CFT, and it does
not fix $\Lambda$.

\begin{admission}
I derived the cosmological minus sign from Schwarzschild--de~Sitter
and recorded the local implication that was already isolated. I did
not produce a holographic pair, a net of CHM balls in de~Sitter, or
a reason to award a prize. The vacuum is de~Sitter. The first law
that selected Einstein$+\Lambda$ is still local Clausius plus an
area law, not FGHMV in cosmology.
\end{admission}

\begin{thebibliography}{99}
\bibitem{Jacobson1995}
T.~Jacobson, Phys.\ Rev.\ Lett.\ 75 (1995) 1260, arXiv:gr-qc/9504004.
\bibitem{GibbonsHawking1977}
G.W.~Gibbons and S.W.~Hawking, Phys.\ Rev.\ D 15 (1977) 2738.
\bibitem{Banihashemi2023}
B.~Banihashemi, T.~Jacobson, A.~Svesko and M.~Visser,
JHEP 01 (2023) 054, arXiv:2208.11706.
\bibitem{McGwierXVIII}
R.W.~McGwier,
``Jacobson is Not a [P] Substitute,''
\url{https://github.com/n4hy/New_Model_Emergent_Gravity}.
\end{thebibliography}

\end{document}
